\chapter{Introduction to the Sociology of India}

\section{The Structure of Paper 2}

Paper 2 (Indian Society) is divided into three sections: Section A, Section B and Section C. Paper 1, in contrast, has no sections --- Sociology is the only optional subject in which Paper 1 carries no sectioning, while Paper 2 is sectioned.

\begin{KeyConcept}
\begin{itemize}
\item \textbf{Section A --- how to study India:} through Western theories or through the Indian (Brahmanical/Indological) perspective.
\item \textbf{Section B --- what to study in India:} caste, village, tribe and kinship (the specifically Indian phenomena).
\item \textbf{Section C --- social change in India:} the most important section in terms of contemporary events.
\end{itemize}
\end{KeyConcept}

Every unit of Section C has one thing in common: all of them are about social change in India --- urbanisation, industrialisation, globalisation, secularism and their impact on family, caste, village and tribe. Section A's theme is the method (how to study India); Section B's theme is the subject matter (caste, village, tribe, kinship).

\begin{DataFact}
Contemporary questions in Paper 2 range between 40 and 70 marks (about 50 marks on average); the theoretical part is roughly 180--200 marks.
\end{DataFact}

\begin{CriticalPoint}
Paper 2 is \textbf{not} dynamic while Paper 1 is static --- both are equally static and equally dynamic. The myth arises because people do not actually know the thinkers (they have read Ghurye but do not know Ghurye).
\end{CriticalPoint}

\begin{QuickRevision}
\begin{itemize}
\item Paper 2 has three sections (A, B, C); Paper 1 has none.
\item Section A = how to study India; Section B = caste/village/tribe/kinship; Section C = social change.
\item Both papers are equally static and dynamic (myth to be avoided).
\end{itemize}
\end{QuickRevision}

\section{Studying India: Western Theories vs Indian Perspective}

There are two broad approaches to studying India: through \textbf{Western theories}, or through the \textbf{Indian perspective}. The teacher deliberately writes ``Indian perspective'' and not ``Indian theories'' --- India, in sociology specifically, has no theories of its own.

\begin{CriticalPoint}
Most Indian sociologists are ``Western by mind and Indian by thinking''; their minds are colonised by Western theory. Talk to any Indian sociologist and they will speak of structural functionalism, feminism or functionalism --- but never of an Indian theory. India could only develop a \textbf{perspective}, not a theory.
\end{CriticalPoint}

The Western theories through which India has been studied are: \textbf{Structural Functionalism}, \textbf{Marxism}, \textbf{Subaltern}, and \textbf{Feminism}. The Indian perspective is \textbf{Indology (Orientalism)}.

\begin{KeyConcept}
``Indian perspective'' does not mean ``study by Indians alone.'' A Britisher or a German scholar can also study India through the Indian perspective. You do not have to be Indian to study India.
\end{KeyConcept}

\begin{ExampleBox}
\textbf{Cultural relativism} --- the teacher's example: in London (2010), a girl of about ten came up and asked him for a lighter; he was shocked and embarrassed, for a ten-year-old smoking would be unimaginable in India, yet in London it is entirely normal. This cultural shock is \textbf{cultural relativism}: every culture is unique, and one is in no position to judge another through one's own culture. Hence studying India through Marxism (a Western product) is studying it through a Western lens.
\end{ExampleBox}

\begin{QuickRevision}
\begin{itemize}
\item Two approaches: Western theories vs Indian perspective.
\item India has no theory, only a perspective (minds colonised by the West).
\item Western theories: Structural Functionalism, Marxism, Subaltern, Feminism; Indian: Indology.
\item Cultural relativism: don't judge another culture through your own.
\end{itemize}
\end{QuickRevision}

\section{Structural Functionalism}

The two most important scholars for structural functionalism are \textbf{M. N. Srinivas} (given in the syllabus) and \textbf{S. C. Dube} (not in the syllabus, but important --- UPSC has already asked a question on him).

\begin{CriticalPoint}
Paper 1 lists six thinkers explicitly, yet questions are routinely asked on thinkers whose names are not listed. There is nothing ``out of syllabus'' in UPSC --- we only know previous-year questions and a few bullet points that we call a syllabus.
\end{CriticalPoint}

\begin{QuickRevision}
\begin{itemize}
\item Structural functionalism: M. N. Srinivas (in syllabus) and S. C. Dube (not in syllabus, but asked).
\item There is nothing ``out of syllabus'' in UPSC.
\end{itemize}
\end{QuickRevision}

\section{Marxism}

The important Marxist scholars who studied India are:
\begin{itemize}
\item \textbf{Dange} --- founder of the Communist Party of India; an academician regarded as one of the first Indian social scientists to study India through the Marxist perspective.
\item \textbf{D. D. Kosambi} --- an important Marxist scholar.
\item \textbf{D. P. Mukerji} --- spelled M-U-K-E-R-J-I (pronounced and written as Mukerji).
\item \textbf{Ramakrishna Mukherjee}.
\item \textbf{A. R. Desai} --- given in the syllabus (named alongside the Marxist perspective).
\end{itemize}

\begin{QuickRevision}
\begin{itemize}
\item Marxist scholars: Dange, D. D. Kosambi, D. P. Mukerji, Ramakrishna Mukherjee, A. R. Desai.
\item A. R. Desai is the one named in the syllabus.
\end{itemize}
\end{QuickRevision}

\section{The Subaltern Perspective}

The subaltern approach is missing from the syllabus, yet UPSC has asked questions on Ambedkar and Ranajit Guha. The important subaltern scholars are:
\begin{itemize}
\item \textbf{B. R. Ambedkar} --- dealt with the untouchables.
\item \textbf{David Hardiman} --- dealt with tribals.
\item \textbf{Ranajit Guha} --- dealt with tribals.
\item \textbf{Gail Omvedt} --- an American sociologist who came to India, married an Indian, stayed, and (the teacher notes) died about a week before this lecture; studied India through the subaltern lens, analysed Dalits and Dalit women, and took three--four factors into account.
\end{itemize}

\begin{PYQLink}
Subaltern is not in the syllabus, but Ambedkar and Ranajit Guha have been asked; expect a question on Gail Omvedt. Many students do not even realise Omvedt is a theorist (and a woman) from the subaltern tradition.
\end{PYQLink}

\subsection{What does ``subaltern'' mean?}

The literal meaning comes from ``sub'' (below), as in ``subway'' --- a way that passes beneath the road. Subaltern, therefore, means looking at India from below --- from the perspective of the most oppressed.

\begin{ComparisonBox}
\begin{itemize}
\item \textbf{Marxism} talks about the exploitation of the \textbf{poor} --- the working class (industrial ``have-nots'') and the peasant/serf (feudal ``have-nots'') --- and this exploitation is purely \textbf{economic}.
\item \textbf{Subaltern} also deals with the poor, but recognises that ``poor'' is not a homogeneous group. Different social locations produce different problems, so it also talks about women and, within them, Dalit women.
\end{itemize}
\end{ComparisonBox}

A Muslim woman and a poor man may both be poor, but the Muslim woman faces a specific patriarchy operating within Muslim households that the poor man does not. Through the intersectionality approach, no class is homogeneous. Because Marxism speaks only of the poor, it cannot address women who are not merely poor but also women --- exploited not only by rich men but by upper castes.

\begin{KeyConcept}
\textbf{Who are the subaltern?} Not women, but \textbf{Dalit women}; not all castes, but the most oppressed castes --- the \textbf{Shudras or the untouchables}.
\end{KeyConcept}

\begin{itemize}
\item \textbf{Shudras and untouchables are different} --- they are not the same; they are grouped here only for convenience.
\item The \textbf{untouchables} are victims of poverty \emph{and} of pollution --- two different things with different causes. Marxism can explain economic exploitation, but not social exploitation or historical injustice.
\item \textbf{Tribals} cannot be called a working class: their issues are forest conservation and displacement-induced development (metro, corridors, railways, airports, highways --- all resting on deforestation and mining). A tribal is not even regarded as Hindu --- seen as outside the Hindu fold.
\end{itemize}

\begin{ExampleBox}
When we celebrate ``20\% growth,'' what does that growth mean for the Dalit woman who cannot go to work, who faces abuse, or for the poor young men with no job? To understand India from the subaltern view is to ask what GDP means to the highly oppressed.
\end{ExampleBox}

\begin{QuickRevision}
\begin{itemize}
\item Subaltern = looking at India from below (sub = below, as in subway).
\item Unlike Marxism (which is economic only), subaltern recognises multiple social locations (intersectionality).
\item Names: Ambedkar (untouchables), Hardiman \& Ranajit Guha (tribals), Gail Omvedt (Dalits/Dalit women).
\item Subaltern groups: Dalit women, Shudras, untouchables, tribals.
\end{itemize}
\end{QuickRevision}

\section{The Feminist Perspective}

The important feminist scholars who studied India are:
\begin{itemize}
\item \textbf{Uma Chakravarti} --- academically recognised as a feminist (people also call Raja Ram Mohan Roy or Savitribai Phule feminist, but academically Chakravarti is the feminist). She described the evolution of \textbf{Brahminical Patriarchy} --- a very different kind of patriarchy.
\item \textbf{Leela Dube} --- wife of S. C. Dube (a structural functionalist), but she herself falls under the feminist perspective.
\item \textbf{Jyotiba Phule and Savitribai Phule} --- largely studied Dalits and Dalit women.
\item \textbf{Sharmila Rege} --- studied Dalit women specifically.
\end{itemize}

\begin{KeyConcept}
\textbf{Brahminical Patriarchy (Uma Chakravarti):} a specific Indian form of patriarchy. After every crime against a Dalit woman, newspapers carry the slogan ``Smash Brahminical Patriarchy'' --- but few know what it means.
\end{KeyConcept}

\subsection{Women as Gateway to Caste (Leela Dube)}

Leela Dube (a follower of Chakravarti) holds that the exploitation of women in India is largely the result of \textbf{caste as a structure}. She also calls women the \textbf{gateway to caste} --- because it is only through women that the caste system reproduces itself.

\begin{ComparisonBox}
\begin{itemize}
\item \textbf{Anuloma (hypergamy):} an upper-caste man marrying a lower-caste woman --- permissible.
\item \textbf{Pratiloma (hypogamy):} a lower-caste man marrying an upper-caste woman --- not permissible (traditionally unacceptable).
\end{itemize}
\end{ComparisonBox}

\begin{ExampleBox}
Love marriage was ``discovered'' in the West (marriage implicitly meant love marriage); arranged marriage was ``made in India'' --- the British discovered it as an Indian institution and had to name it. Arranged marriage means marriage \textbf{within the same caste}: the two people may differ, but the caste is already fixed, because the family has socialised the girl into that caste.
\end{ExampleBox}

Girls are preserved by the family as pure, and their wishes are sidelined, so that through this girl the family's caste is reproduced in the next generation. The girl is exchanged between two similar castes --- hence she is the gateway to caste.

\subsection{Structuralism (a new theoretical perspective)}

\begin{KeyConcept}
\textbf{Structuralism} is all about language: communication is possible only through the exchange of words. Applied to marriage, a marriage is nothing but a communication facilitated not by the exchange of words but by the \textbf{exchange of girls} --- girls are exchanged so that two families can communicate and form a new alliance. The girl becomes an object circulated from one generation to the next.
\end{KeyConcept}

\begin{QuickRevision}
\begin{itemize}
\item Feminist scholars: Uma Chakravarti (Brahminical Patriarchy), Leela Dube (women as gateway to caste), Phules, Sharmila Rege.
\item Anuloma = hypergamy (upper man + lower woman, permissible); Pratiloma = hypogamy (prohibited).
\item Structuralism: marriage = exchange of girls (like exchange of words).
\end{itemize}
\end{QuickRevision}

\section{The Indian Perspective: Indology \& Orientalism}

Whatever we know of India today is largely \textbf{colonial knowledge}. The guiding question: do we know India through ourselves, or through the British? The answer the teacher gives: we know India through a \textbf{Brahminical lens}.

\begin{KeyConcept}
\textbf{Indology} simply means the study of India by anyone --- but with a specific method. Studying India through the ancient Brahmanical texts is the \textbf{Indological approach}.
\end{KeyConcept}

Imagine being a British professor teaching British students: you introduce Indology (India), Egyptology (Egypt), Tibetology (Tibet/China) --- all ``Asian cultures.'' This is Orientalism. The West ``invented man,'' and this invented man ``discovered'' India as the ``other'' culture.

\begin{ExampleBox}
A traveller hires a guide; after daily conversation a bond forms, and the traveller comes to accept whatever the guide says as reality. The British did the same: whatever the \textbf{Brahmins} told them became ``reality.'' Hence we know India through the British, but the British knew India through the Brahmins --- a \textbf{Brahminical lens}.
\end{ExampleBox}

The Brahmins introduced the British to the traditional texts --- the \textbf{Dharma Shastras, Dharma Sutras, Smritis and Manu Smriti} --- which were written in Sanskrit. Learning Sanskrit became the British's first task, because once they could read Sanskrit they no longer needed the guide.

The result: studying India through the ancient Brahmanical texts is the Indological approach --- a \textbf{book view} of India.

\subsection{The founders of Indology}

\begin{itemize}
\item \textbf{William Jones} --- the first to study India through ancient Brahmanical texts in a systematic form; established the \textbf{Asiatic Society} (c.~1774/1775--80s) and its journal \textbf{Asiatic Researches}; translated the \textbf{Rig Veda}; published \textbf{Sacred Books of the East}; a contemporary of Warren Hastings (who invited him --- Jones was a Supreme Court judge) and of Wellesley.
\item \textbf{Wilkins} --- translated the \textbf{Bhagavad Gita} and the \textbf{Upanishads}; first learned Sanskrit.
\item \textbf{Max M\"uller} --- a German; gave the theory of \textbf{naturism} (nature worship as the oldest religion); the first to see similarity between Sanskrit and German; professor of \textbf{comparative philology} at Oxford; coined \textbf{Aryan Brotherhood}.
\end{itemize}

\begin{DataFact}
\textbf{A. L. Basham} (author of ``The Wonder That Was India'') called William Jones and Wilkins the \textbf{true founders of Indology in India}.
\end{DataFact}

The first similarity the British found between India and Europe was in \textbf{language} --- Sanskrit resembles European languages, especially German. William Jones went so far as to claim that all European languages have their origin in Sanskrit. This gave the world the word \textbf{Indo-European} (from comparative philology).

\begin{KeyConcept}
Four beliefs common to Jones, Wilkins and Max M\"uller: (1) Indian culture is Brahminical culture; (2) to understand India one must learn Sanskrit; (3) Indians and Europeans are similar; (4) Sanskrit is the origin of the other (European) languages.
\end{KeyConcept}

William Jones established a \textbf{Department of Sanskrit and Indic Studies} within the Asiatic institution, thereby making Indology a discipline in itself --- like physics or chemistry (just as we now have South Asian studies at JNU and a South Asian University in Delhi).

\begin{QuickRevision}
\begin{itemize}
\item Indian perspective = Indology (book view via ancient Brahmanical texts).
\item Founders: William Jones (Asiatic Society, Rig Veda, Sacred Books of the East), Wilkins (Gita, Upanishads), Max M\"uller (Aryan Brotherhood).
\item We know India through a Brahminical lens.
\end{itemize}
\end{QuickRevision}

\section{How We Know India: Travellers, Missionaries, Orientalists}

\textbf{Bernard Cohn} (a Chicago anthropologist) describes three routes through which we know India: travellers' accounts, missionaries' accounts, and the accounts of the Orientalists.

\subsection{Travellers}

\begin{itemize}
\item \textbf{Megasthenes} --- wrote \textbf{Indica}; visited India during Chandragupta Maurya's time, sent by the Seleucid Greek Empire; noted seven groups in the Mauryan administration. He observed the \textbf{caste system but not slavery} (though Kautilya's Arthashastra does record slavery --- a reminder of how ambiguous the early accounts, and the very word ``caste,'' are); he had a \textbf{field view} of India.
\item \textbf{Ibn Battuta} --- wrote the \textbf{Rihla}; knew Arabic, Sanskrit and Greek; learned Sanskrit quickly and befriended Brahmins; witnessed and wrote on the caste system.
\item \textbf{Abdul Razak Samarkandi} --- well versed in Sanskrit, also wrote on India.
\item \textbf{Fran\c{c}ois Bernier} --- visited during the time of Dara Shikoh; \textbf{Tavernier} (a Frenchman) also came.
\end{itemize}

\subsection{Missionaries}

Missionaries came to convert (proselytise) people into Christianity. Because their motive was conversion, their account of other religions is \textbf{value-loaded} --- they treat other religions as inferior.

\subsection{Orientalists}

\begin{KeyConcept}
\textbf{Orientalism} is the study of the ``Orient'' (the Eastern/other cultures) by Westerners. A Westerner explaining India is an Orientalist --- and an Orientalist is always an outsider; an Indian cannot be an Orientalist. For convenience, treat \textbf{Indologist and Orientalist as synonymous} (the difference only arises with Ghurye's Indian Indology).
\end{KeyConcept}

\begin{QuickRevision}
\begin{itemize}
\item Bernard Cohn's three routes: travellers, missionaries, orientalists.
\item Travellers: Megasthenes (Indica), Ibn Battuta (Rihla), Samarkandi, Bernier, Tavernier.
\item Orientalist = always an outsider; Indologist treated as synonymous with Orientalist (for now).
\end{itemize}
\end{QuickRevision}

\section{Book View vs Field View}

\begin{ComparisonBox}
\begin{itemize}
\item \textbf{Book view:} understanding India through ancient texts (Vedas, Upanishads, epics) --- the method of William Jones and Wilkins.
\item \textbf{Field view:} understanding India by going into the field, living in villages, observing directly --- the method of the travellers and later of field-based sociology.
\end{itemize}
\end{ComparisonBox}

Sociology emerged in the West in response to urbanisation and industrialisation --- not to rural settings. Studying villages is therefore a uniquely Indian exercise. \textbf{Andre Beteille} notes that Indian sociologists are regarded abroad as \textbf{social anthropologists}, even while they practise sociology at home.

\begin{CriticalPoint}
A hierarchy was created between the two disciplines: sociology was for the modern, civilised, Western, educated people; anthropology was for the ``other,'' uncivilised, tribal cultures. Today these distinctions have lost relevance --- anthropologists study the West as much as sociologists study the East.
\end{CriticalPoint}

In its early days, Indian sociology was called \textbf{social anthropology}. Ghurye was an anthropologist and Indologist; his students (Iravati Karve, M. N. Srinivas) were likewise regarded as social anthropologists.

\begin{QuickRevision}
\begin{itemize}
\item Book view (texts) vs Field view (living in villages).
\item Early Indian sociology = social anthropology (hierarchy: sociology > anthropology).
\item Ghurye and his students (Iravati Karve, M. N. Srinivas) were social anthropologists.
\end{itemize}
\end{QuickRevision}

\section{Section B: Village, Tribe, Caste, Kinship}

Four phenomena are seen as particularly Indian: village, tribe, caste and kinship.

\subsection{Tribe}

\textbf{G. S. Ghurye} was the first to open the debate on tribe in Indian sociology, in debate with the British anthropologist \textbf{Verrier Elwin}.

\begin{ComparisonBox}
\begin{itemize}
\item \textbf{Ghurye (nationalist):} tribals are Hindus --- merely \textbf{backward Hindus} who, if assimilated, will gradually become forward Hindus. Isolating them traps them in a permanent cycle of isolation.
\item \textbf{Elwin (isolationist):} tribal cultures must be preserved in isolation. This view won --- including in the Constituent Assembly.
\end{itemize}
\end{ComparisonBox}

\begin{GovtPolicy}
The isolationist view survives today in the Fifth and Sixth Schedules and the \textbf{Inner Line Permit} (North-East) that protect tribal pockets. The critique: tribals have become ``artifacts in a museum,'' looked at like animals.
\end{GovtPolicy}

\subsection{Caste}

\begin{KeyConcept}
Indian sociology is nothing less than the \textbf{sociology of caste} --- the two have become synonymous. Indian sociologists remain focused on caste (seen as traditional), not on class (individualism and meritocracy).
\end{KeyConcept}

\begin{CriticalPoint}
The caste paradox: to end caste you must first know what caste is; but the moment you explain caste to someone (a father to his child), caste gets a new lease of life and reproduces itself. The dream of caste being dead dies with every explanation.
\end{CriticalPoint}

\subsection{Kinship}

Families are formed around the notion of caste. If 93\% of India practises endogamy, then 93\% of families are formed around caste --- so kinship networks are caste-based. In rural India there are even caste-based moneylender associations.

\begin{DataFact}
\textbf{93\% of India practises endogamy}, so roughly 93\% of Indian families are formed around caste.
\end{DataFact}

\begin{KeyConcept}
Caste is not just identity --- it is \textbf{capital} and a \textbf{resource}: it forms family, kinship, alliances and vote banks, and is used for mobility (reservation) and political power. ``To say caste should end is to say life should end in India.''
\end{KeyConcept}

\textbf{David Mandelbaum} identified \textbf{village, caste, kinship and tribe} as the four basic units of Indian society.

\begin{QuickRevision}
\begin{itemize}
\item Section B's four Indian phenomena: village, tribe, caste, kinship.
\item Ghurye (assimilate) vs Elwin (isolate) --- isolation won.
\item Caste = resource/capital; 93\% endogamy; Mandelbaum's four basic units.
\end{itemize}
\end{QuickRevision}

\section{Section C: Social Change in India}

Section C is the most important in terms of contemporary events. Its unifying theme is social change in India.

\subsection{Village}

\begin{itemize}
\item Villages are being urbanised; urban development schemes reach rural areas every year --- this is urbanisation.
\item \textbf{Bharat Nirman} means the reconstruction of villages. Gandhi believed India lies in the village (``Bharat lies in village''), so village and India are synonymous.
\end{itemize}

\subsection{Tribe}

\begin{itemize}
\item Since independence many tribals have been assimilated into a non-tribal culture (Hindu or Western). The \textbf{7.5\% reservation for STs} is a ``gateway to become a member of another culture'' they never belonged to.
\item Tribals who resisted the British (e.g.\ the \textbf{Tana Bhagat movement}) are today becoming doctors, engineers and IES officers --- accommodated into the meritocratic, Western fold.
\item The \textbf{Criminal Tribes} policy was done away with; many such tribes are now classified as Scheduled Tribes and given reservation --- the moment they are given reservation, they are modernised.
\end{itemize}

\subsection{Kinship / family}

\begin{itemize}
\item The Western theory holds that the family life-cycle moves from jointness to nuclearity (the joint family breaking into nuclear households).
\item The idea of the Indian joint family is known to us largely through the British --- via the epics (the Pandavas, the Kauravas, five or more brothers under one roof).
\item Today, jobs in Bangalore, Delhi or America create \textbf{migrant households, not families}; the joint family (with joint residence) is dissolving.
\end{itemize}

\subsection{Secularism \& globalisation}

Secularism is a Western idea that altered Indian power structures; globalisation and its impact are read on family, elderly, caste and politics --- always social change.

\begin{QuickRevision}
\begin{itemize}
\item Section C = social change: urbanisation, Bharat Nirman, 7.5\% ST reservation, criminal tribes reclassified as ST.
\item Joint family turning nuclear; migrant households; secularism; globalisation.
\end{itemize}
\end{QuickRevision}

\section{Misconceptions About Paper 2}

\begin{CriticalPoint}
The claim ``Paper 1 is static, Paper 2 is dynamic'' is a wrong classification. Both are equally static and equally dynamic. People say this because they confuse the ~50 marks of contemporary questions with the full 250 marks.
\end{CriticalPoint}

The teacher closes by flagging that \textbf{caste} and \textbf{Varna} are often used interchangeably --- M. N. Srinivas wrote an article titled \textbf{``Varna and Caste''} precisely to separate the two. This is taken up the next day.

\begin{QuickRevision}
\begin{itemize}
\item Paper 2 is not dynamic-only --- the static/dynamic split is a myth.
\item Caste vs Varna: M. N. Srinivas, ``Varna and Caste.''
\end{itemize}
\end{QuickRevision}

\section{Preview: Louis Dumont}

\textbf{Louis Dumont} is an Indologist --- but unlike Ghurye he did not rely only on Brahmanical texts. He also did field work: in Tamil Nadu (the \textbf{Pramalai Kallar} caste of Madurai district, at \textbf{Tengal Patti} --- once a criminal tribe) and in \textbf{Gorakhpur}. From that field work he derived his theory of \textbf{purity and impurity}.

\begin{ComparisonBox}
\begin{itemize}
\item \textbf{G. S. Ghurye} is mainly \textbf{book view}.
\item \textbf{Louis Dumont} is \textbf{book view + field view}.
\end{itemize}
\end{ComparisonBox}

\begin{QuickRevision}
\begin{itemize}
\item Louis Dumont = Indologist with book + field view; purity/impurity theory (Tamil Nadu, Gorakhpur).
\item Ghurye = book view only.
\end{itemize}
\end{QuickRevision}
